\section{附带民事诉讼}
\begin{enumerate}
  \item 核心。
        \begin{enumerate}
          \item \bs{必须以刑事诉讼成立为前提；刑主民从}。
          \item \bs{犯罪行为造成的实际的、必然的、物质损坏}。
          \item \bs{本质上是个民诉}。\zd{刑事若无罪，附民可调、可判、可另诉}。
        \end{enumerate}
  \item 当事人。因犯罪行为遭受物质损坏。
        \begin{enumerate}
          \item 原告。\zd{被害人自己}。死了 \yjt 近亲属 \hz 法定代理人。\al{路人甲受害}。
          \item 国家、集体遭到了损失，检察院也可做原告人。
          \item 被告人 \bdy 刑事被告人。
          \item 立案后及时提，二审提只给调。\zd{调解协议没有强制执行力}。履行完毕有终结效力，无法反悔。
          \item 赔偿范围。\bs{医护交康加误工，残疾辅助死亡葬}。
        \end{enumerate}
  \item 法院。
        \begin{enumerate}
          \item 公诉撤回起诉的，附民调解不成裁定驳回，告知另诉。
          \item 共同侵害人，原告人有选择权。只告一部分\zd{应当告知}有权对其他人一并起诉。例外。共犯在逃、已从其他同案犯得到足额赔偿。
          \item 不予支持。
                \begin{enumerate}
                  \item \zd{不赔精神损害}。\yi 有其他的途径来安抚他的精神创伤。\er \zd{判了空判}。死刑犯无所畏惧。
                  \item \zd{非法占有、处置的财产犯罪。追缴、退赔}。\al{被变卖的手机}。参看涉案财物处理。
                  \item 通过附民要求国赔的，\bs{只能单独申请国赔}。
                \end{enumerate}
          \item 审理。类推民诉。
                \begin{enumerate}
                  \item 刑和民\zd{同一个审判组织}；如果分开审，\zd{先刑后民}；民事和刑事\zd{不能相抵触}。
                  \item 民事过分拖延，不影响刑事裁判的正常生效。
                  \item \zd{不收诉讼费}。
                  \item 原告不来 \hz 中途退，撤诉处理；共监继不来，可以缺席。
                \end{enumerate}
          \item 附带民诉的财产保全问题。\bs{审理刑事案件的法院采取保全措施}。其他同民诉。
          \item 二审附民。
                \begin{enumerate}
                  \item 上诉期跟着刑诉来。
                  \item \zd{可暂缓送监执行}。
                  \item 增加独立请求 \hz 提起反诉只能调解。
                  \item 分开生效。
                        \begin{enumerate}
                          \item 刑上民不上。刑\sjt/民 \yjt 生效 \yjt 有错 \yjt 民再审。
                          \item 民上刑不上。民\sjt/刑 \yjt 生效 \yjt 有错 \yjt 刑再审 \yjt 带上民。刑事错误会连累民诉。
                        \end{enumerate}
                \end{enumerate}
          \item 刑事附带民事公益诉讼，由\zd{审理刑事案件的法院管辖}。
        \end{enumerate}
\end{enumerate}
